\section{Comparação A/B com estudantes}
\label{sec:avaliacao-estudantes}

Esta avaliação foi feita com uma turma de estudantes, que comparou duas versões do mesmo programa: a Versão A, na sintaxe padrão, e a Versão B, personalizada em português. O formulário também traz afirmações sobre a plataforma e a demonstração.

\subsection{Formulário e procedimento}
\label{subsec:formulario-estudantes}

O formulário dos estudantes não constava da versão anterior do protocolo. Ele foi aplicado em 3 de novembro de 2026 a uma turma do 1º ano do ensino médio técnico em Informática da UFCI-MG, Campus Leopoldina. Os autores demonstraram a plataforma em sala, os alunos também a usaram e responderam em seguida, entre 14h13 e 14h43. Dezoito estudantes responderam ao formulário.

O formulário é curto. Depois do termo de assentimento, pergunta a idade, o curso, o ano e se o aluno já tinha programado. Em seguida traz seis afirmações em escala de cinco pontos, de ``discordo totalmente'' a ``concordo totalmente'':

\begin{itemize}
    \item entendi para que serve a plataforma;
    \item um programa escrito com palavras em português é mais fácil de entender do que um escrito com palavras em inglês;
    \item poder escolher as palavras da linguagem ajudaria quem está começando a programar;
    \item ver como a plataforma separa o programa em partes, os \textit{tokens}, ajuda a entender como o computador lê o código;
    \item eu gostaria de usar essa plataforma nas aulas de programação;
    \item a demonstração aumentou meu interesse em aprender programação.
\end{itemize}
 Depois o formulário mostra duas versões de um mesmo programa, que lê um número e escreve se ele é par ou ímpar. A Versão A usa a sintaxe padrão, em inglês, e a Versão B usa uma sintaxe personalizada, em português. O aluno indica qual das duas acha mais fácil de entender, ou se considera as duas igualmente fáceis ou igualmente difíceis. A medida obtida com essa pergunta é a preferência declarada entre as versões, isto é, o número de alunos que escolheu cada uma das quatro alternativas; ela registra o que o aluno diz achar mais fácil e não mede compreensão. As duas versões têm as mesmas instruções, na mesma ordem e com a mesma disposição em linhas. Mudam as palavras reservadas, o operador de igualdade, que na Versão B é escrito \texttt{igual}, e os delimitadores de bloco, que passam das chaves para \texttt{inicio} e \texttt{fim}. Elas aparecem como texto simples no corpo da pergunta, sem imagem nem realce de cores, e estão reproduzidas no Apêndice~\ref{ap:formulario-estudantes}. O último campo é um comentário livre sobre o que mudaria ou acrescentaria na plataforma.

\subsection{Resultados da sessão}
\label{sec:resultado-estudantes}

Esta subseção trata das afirmações sobre a demonstração (questões \ref{qF:demo-ini} a \ref{qF:demo-fim}) e da pergunta de comparação entre as versões (questão \ref{qF:comparacao}) do Apêndice~\ref{ap:formulario-estudantes}. O formulário descrito na Subseção~\ref{subsec:formulario-estudantes} foi aplicado em 3 de novembro de 2026 a uma turma do 1º ano do ensino médio técnico em Informática da UFCI-MG, Campus Leopoldina, depois de uma demonstração da plataforma feita pelos autores e de um período em que os alunos também a usaram. Houve dezoito respostas, todas com o termo de assentimento marcado. Nove alunos tinham até 15 anos e nove tinham 16 ou 17. Nove disseram já ter programado um pouco, oito disseram programar com frequência e um disse nunca ter programado. Quase ninguém, portanto, era iniciante absoluto.

A Tabela~\ref{tab:resultado-estudantes} traz as seis afirmações. Todas tiveram média entre 4,00 e 4,33.

\begin{table}[htbp]
\centering
\caption{Afirmações do formulário da sessão com estudantes (escala de 1 a 5; dezoito respostas).}
\label{tab:resultado-estudantes}
\footnotesize
\renewcommand{\arraystretch}{1.2}
\begin{tabular}{|p{6.3cm}|c|c|c|c|c|}
\hline
\textbf{Afirmação} & \textbf{Média} & \textbf{Mediana} & \textbf{Concordam} & \textbf{Neutros} & \textbf{Discordam} \\\hline
Entendi para que serve a plataforma & 4,33 & 4,5 & 17 & 0 & 1 \\\hline
Um programa escrito com palavras em português é mais fácil de entender do que um escrito com palavras em inglês & 4,00 & 4 & 11 & 6 & 1 \\\hline
Poder escolher as palavras da linguagem ajudaria quem está começando a programar & 4,17 & 4 & 14 & 3 & 1 \\\hline
Ver como a plataforma separa o programa em partes (os \textit{tokens}) ajuda a entender como o computador lê o código & 4,17 & 4 & 15 & 1 & 2 \\\hline
Eu gostaria de usar essa plataforma nas aulas de programação & 4,17 & 4 & 15 & 3 & 0 \\\hline
A demonstração aumentou meu interesse em aprender programação & 4,17 & 4 & 16 & 0 & 2 \\\hline
\end{tabular}
\legend{Fonte: dados da pesquisa (2026).}
\end{table}

Dezessete dos dezoito alunos disseram ter entendido para que serve a plataforma, e quinze disseram que gostariam de usá-la nas aulas; nesse item ninguém discordou. A afirmação de que um programa em português é mais fácil de entender do que um em inglês teve a menor média e o maior número de neutros: onze concordaram, seis ficaram no ponto neutro e um discordou. Esse é o item mais próximo da premissa da plataforma, e um terço da turma não tomou posição sobre ele.

O formulário também mostrou duas versões de um programa que lê um número e escreve se ele é par ou ímpar: a Versão A, na sintaxe padrão em inglês, e a Versão B, personalizada em português. As duas têm as mesmas instruções e a mesma disposição e diferem nas palavras reservadas, no operador de igualdade e nos delimitadores de bloco; o texto de ambas está no Apêndice~\ref{ap:formulario-estudantes}. A Tabela~\ref{tab:resultado-versoes} separa as respostas pela experiência declarada. Onze alunos acharam a Versão B mais fácil de entender, quatro acharam as duas igualmente fáceis, dois preferiram a Versão A e um achou as duas igualmente difíceis.

\begin{table}[htbp]
\centering
\caption{Versão do programa considerada mais fácil de entender, por experiência declarada com programação.}
\label{tab:resultado-versoes}
\footnotesize
\renewcommand{\arraystretch}{1.2}
\begin{tabular}{|p{4.3cm}|c|c|c|c|c|}
\hline
\textbf{Experiência} & \textbf{Versão B} & \textbf{Igualmente fáceis} & \textbf{Versão A} & \textbf{Igualmente difíceis} & \textbf{Total} \\\hline
Programa com frequência & 6 & 2 & 0 & 0 & 8 \\\hline
Já programou um pouco & 5 & 2 & 2 & 0 & 9 \\\hline
Nunca programou & 0 & 0 & 0 & 1 & 1 \\\hline
\textbf{Total} & 11 & 4 & 2 & 1 & 18 \\\hline
\end{tabular}
\legend{Fonte: dados da pesquisa (2026).}
\end{table}

A preferência pela versão em português apareceu nos dois grupos com experiência, e os dois alunos que preferiram a versão em inglês estão entre os que já tinham programado um pouco. O único aluno que nunca tinha programado, justamente o perfil para o qual a plataforma foi pensada, achou as duas versões igualmente difíceis e discordou de três das seis afirmações, entre elas a de ter entendido para que serve a plataforma. Uma única resposta não sustenta conclusão, mas ela vai em sentido contrário ao das demais e fica registrada.

O campo de comentário livre recebeu duas respostas: uma elogiou a plataforma e a outra pediu o acréscimo de outra linguagem. O formulário não tem item de autorização de citação, e por isso os comentários não são transcritos.

Pelos limites descritos na Seção~\ref{subsec:limites-interpretacao}, essas respostas servem como indício de aceitação e de preferência declarada pela versão em português e não mostram redução da barreira sintática nem aprendizagem. O registro do assentimento e a falta de autorização dos responsáveis estão descritos na Seção~\ref{sec:protocolo-avaliacao} e retomados nas limitações do Capítulo~\ref{cap:consideracoes}.

